\documentclass[11pt]{article}
\usepackage{amsmath,amssymb,amsthm,hyperref}
\newtheorem{proposition}{Proposition}
\title{Literature bridges: three conditional reductions}
\author{Internal research audit}
\date{15 November 2026}
\begin{document}
\maketitle
These are applications of named external theorems, not solutions of HC4, JC2,
the exceptional Mathieu statements, or the remaining Wright families.
Independent review is pending. The accompanying JSON records applicability,
prior adoption, source versions and the missing hypotheses.

\begin{proposition}[Fixed-curve moment shift]
Let $R$ be a compact Riemann surface, $f$ a nonconstant meromorphic function,
$\omega$ a meromorphic one-form with poles only at poles of $f$, and $\gamma$
a closed rectifiable cycle avoiding those poles. If
$\int_\gamma f^m\omega=0$ for every $m\geq1$, then
$\int_\gamma\omega=0$.
\end{proposition}
\begin{proof}
Use Gavrilov--Pakovich, Theorem 2.1 of
\href{https://arxiv.org/abs/1107.3029}{arXiv:1107.3029}.
Its finite branch sums for $\omega$ are
$\varphi_i(z)=\sum_j n_{ij}(\omega/df)(f_j^{-1}(z))$.
The integers depend on $f$ and $\gamma$, so the sums for $f\omega$ are
$z\varphi_i(z)$. All nonnegative moments for $f\omega$ vanish by assumption;
the pole hypothesis is preserved. The theorem gives $z\varphi_i=0$, hence
$\varphi_i=0$, and its reverse implication gives every nonnegative moment
for $\omega$ zero, including the zeroth one.
\end{proof}
An actual Haar representation of this kind with mass one would therefore
exclude a pure phase. A Haar density which is itself a period on a moving
family of curves is not such a representation. Relative paths and finite
poles also require separate treatment. No reduction of general Haar laws
to this fixed-curve setting has been established here.

\begin{proposition}[A short forbidden filling window]
Let $J$ be an L-space knot with $2g(J)-1=2A-1$, where $A\geq1$ is an integer.
For a positive integer $p$, put $d=p-(2A-1)$. If $d<1+\sqrt p$, then
$S^3_p(J)$ cannot bound a smooth negative-definite four-manifold with
torsion-free integral first homology. In particular this excludes such
fillings at $p=2A-1$ and $p=2A+1$. If a filling is sharp, it must additionally
satisfy $d\geq\sqrt{3p+1}$.
\end{proposition}
\begin{proof}
Greene, Theorem 1.1 of
\href{https://arxiv.org/abs/1009.1130v2}{arXiv:1009.1130v2}, gives
$2g(J)-1\leq p-\sqrt p-1$, and the sharp version gives
$2g(J)-1\leq p-\sqrt{3p+1}$.
At the two indicated slopes $d=0,2$ respectively; for the second,
$p=2A+1\geq3$ implies $1+\sqrt p>2$.
\end{proof}
The retained Wright all-plus solid companions have the stated genus.
This is a conditional test for their oriented integer surgeries: a solid
three-dimensional cap or the L-space property does not supply the required
four-dimensional filling. No remaining Wright family is closed by this test.

\begin{proposition}[Finite-precision detection over a fixed valuation ring]
Let $V$ be a Henselian valuation ring with separable fraction-field
completion, $X/V$ a scheme of finite presentation, and
$I_\alpha=\{x:v(x)\geq\alpha\}$ for positive $\alpha$ in its value group.
There are a positive integer $N$ and $\delta\geq0$ such that
\[
\operatorname{Im}\bigl(X(V/I_{N\alpha+\delta})\to X(V/I_\alpha)\bigr)
=\operatorname{Im}\bigl(X(V)\to X(V/I_\alpha)\bigr).
\]
Consequently, a prescribed residue point with no exact $V$-lift has no
approximate lift at some finite valuation precision. Compatibility between
hypothetical approximations is not required.
\end{proposition}
\begin{proof}
The displayed equality is Moret-Bailly, Theorem 1.2 of
\href{https://arxiv.org/abs/1106.0984v3}{arXiv:1106.0984v3}.
Any sufficiently precise approximate lift has an exact lift agreeing modulo
$I_\alpha$, hence modulo the maximal ideal. This proves the consequence.
For a DVR it is also a direct consequence of Greenberg's Theorem 1,
\href{https://numdam.org/articles/10.1007/BF02684802/}{IHES 31 (1966), 59--64}.
\end{proof}
For fixed degree $D$, coefficient equations for $\det\operatorname{Hess}h=1$
define such a finite-presentation scheme. Constants may depend on the scheme
and valuation ring. There is no all-degree or effective bound here. Over the
fixed ring $\mathcal O_{\mathbb C_p}$, the same statement gives an absolute
valuation threshold for residue points with no exact lift over that ring;
a DVR approximation of precision $s$ and ramification $e$ has valuation
precision $s/e$ when $v(p)=1$. This is not an identification of all finite-DVR
lifting problems with lifting over $\mathcal O_{\mathbb C_p}$. For the retained
nonplane/collision examples, the repository already supplies a more explicit
linear-in-$e$ obstruction, so this application is supplementary.
\end{document}
